%% kapa: exam=1 round=21 year=2010 period=2 form=A subject=appraisal_law subject_ko=부동산관계법규 seq=6 number=46
%% source: Q-Net 감정평가사 기출문제 (cst00309) artlSeq=5024181 file=extracted/1008/5024181/2010년 제21회 감정평가사 1차시험 기출문제/2교시.hwp
%% conversion: kordoc 4.24.51 -> tools/gen1.py (auto); review=proofread-gpt-sol; flags=none
\begin{question}{46}
국토의 계획 및 이용에 관한 법령상 도시기본계획에 관한 설명으로 옳지 않은 것은?
\begin{choices}
  \item 특별시장․광역시장․시장 또는 군수는 지역여건상 필요하다고 인정되면 인접한 특별시․광역시․시 또는 군의 관할구역 전부 또는 일부를 포함하여 도시기본계획을 수립할 수 있다.
  \item 시장 또는 군수는 도시기본계획을 수립하거나 변경하려면 대통령령으로 정하는 바에 따라 도지사의 승인을 받아야 한다.
  \item 특별시장․광역시장․시장 또는 군수는 도시기본계획 수립시 광역도시계획 및 국가계획의 내용을 반영하여야 한다.
  \item 도시기본계획과 광역도시계획의 내용이 서로 다를 경우 광역도시계획의 내용이 우선한다.
  \item 특별시장․광역시장이 도시기본계획을 수립하거나 변경할 경우 지방도시계획위원회의 심의를 거친 후 관계행정기관의 장과 협의하여야 한다.
\end{choices}
\end{question}
\answer{5}
